% =============================================================================
%  Chapter 3 — Negative Weights and the Checkerboard Artefact
% =============================================================================

\chapter{Negative Weights and the Checkerboard Artefact}
\label{ch:negative_weights}

This chapter documents the diagnosis and resolution of a persistent
\emph{checkerboard} (or ``diamond'') artefact observed in nodal stress fields
projected from Gauss points onto quadratic tetrahedral elements (TET10).  The
analysis traces the root cause to the interaction between negative quadrature
weights and shape function values, and prescribes two complementary fixes.


% ─────────────────────────────────────────────────────────────────────────
\section{Symptom: Element-Level Checkerboard Pattern}
\label{sec:symptom}

When visualising projected stress or strain fields in ParaView on TET10 meshes,
the following pathology was observed:
\begin{itemize}[nosep]
  \item At \emph{vertex} nodes, the projected field values had the
        \emph{opposite sign} of the expected physical values.
  \item At \emph{edge midpoint} nodes, the values were correct.
  \item The pattern repeated \emph{at every element}, forming a regular
        alternating (checkerboard) structure.
  \item On HEX27 meshes, the same analysis pipeline produced smooth, correct
        fields.
\end{itemize}


% ─────────────────────────────────────────────────────────────────────────
\section{Root Cause Analysis}
\label{sec:root_cause}

The projected nodal values were computed via \emph{lumped $L^2$ projection}:
\begin{equation}
  \tilde{\sigma}_I
  = \frac{\displaystyle\sum_{e \in \mathcal{E}_I}
          \sum_{g=1}^{n_g} N_I(\bxi_g)\, w_g\, |J_g|\, \sigma_g}
         {\displaystyle\sum_{e \in \mathcal{E}_I}
          \sum_{g=1}^{n_g} N_I(\bxi_g)\, w_g\, |J_g|},
  \label{eq:lumped_L2}
\end{equation}
where $\mathcal{E}_I$ is the set of elements sharing node $I$, $N_I$ is the
global shape function, and $\sigma_g$ is the Gauss-point value.

For the projection~\eqref{eq:lumped_L2} to produce meaningful results, the
denominator must be \textbf{strictly positive} for every node.  We will show
that this condition \emph{fails} at vertex nodes of TET10 when the Keast
5-point rule is used.


\subsection{The Keast 5-Point Rule Weights}

The Keast degree-3 rule for the tetrahedron has 5~points
(\cref{sec:keast}):
\begin{equation}
  \begin{aligned}
    &\text{Point 1 (centroid):}   & w_1 &= -\frac{2}{15}\,|\calS_3|
      = -\frac{1}{45} \approx -0.02222, \\
    &\text{Points 2--5 (vertices):} & w_g &= +\frac{3}{40}\,|\calS_3|
      = +\frac{1}{80} = 0.01250.
  \end{aligned}
  \label{eq:keast5_detail}
\end{equation}
The centroid weight $w_1$ is \textbf{negative}.


\subsection{TET10 Vertex Shape Functions at Quadrature Points}

The vertex shape function for node $i$ is (cf.\
\cref{eq:shape_vertex_quad}):
\[
  N_i(\bxi) = \lambda_i\,(2\lambda_i - 1).
\]
At the 5 Keast quadrature points:

\paragraph{At the centroid} ($\lambda_i = 1/4$ for all $i$):
\[
  N_i = \tfrac{1}{4}(2 \cdot \tfrac{1}{4} - 1) = \tfrac{1}{4}(-\tfrac{1}{2})
      = -\tfrac{1}{8}.
\]

\paragraph{At the vertex-orbit points:} Each vertex-orbit point has one
``dominant'' barycentric coordinate near $0.5854$ and three ``small'' ones near
$0.1382$.

For the vertex $i$ that is \emph{dominant} at point $g$:
\[
  N_i \approx 0.5854 \times (2 \times 0.5854 - 1) = 0.5854 \times 0.1708
      \approx 0.1000.
\]

For a vertex $i$ that is \emph{not dominant}:
\[
  N_i \approx 0.1382 \times (2 \times 0.1382 - 1) = 0.1382 \times (-0.7236)
      \approx -0.1000.
\]


\subsection{The Denominator Catastrophe}

Consider the denominator of~\eqref{eq:lumped_L2} for a \emph{vertex} node $I$,
restricted to a single element:
\begin{equation}
  D_I = \sum_{g=1}^{5} N_I(\bxi_g)\, w_g\, |J_g|.
  \label{eq:denominator_element}
\end{equation}

Assuming approximately constant $|J_g| \approx |J|$ within the element:
\begin{align}
  D_I &\approx |J| \bigl[
    \underbrace{N_I(\bxi_1)\,w_1}_{\text{centroid: } (-1/8)(-1/45) > 0}
    +\, N_I(\bxi_2)\,w_2
    + \cdots
    + N_I(\bxi_5)\,w_5
  \bigr].
  \label{eq:D_I_expansion}
\end{align}

The centroid contribution is \emph{positive} (negative times negative), but the
vertex-orbit contributions produce a mix of positive and negative terms.  The
crucial issue is that at the three vertex-orbit points where $N_I < 0$
(because $\lambda_I < 1/2$), the \emph{positive} weight $w_g > 0$ creates a
negative contribution, while at the one vertex-orbit point where $N_I > 0$
(the ``own'' vertex direction), the positive weight creates a positive
contribution.

When these are summed, the negative contributions can \emph{overwhelm} the
positive ones, making $D_I < 0$ or $D_I \approx 0$.


\subsection{Consequence: Sign Flip or Division by Zero}

\begin{equation}
  D_I \le 0 \implies \tilde{\sigma}_I
  = \frac{\text{(numerator)}}{D_I}
  \begin{cases}
    \text{has wrong sign}  & \text{if } D_I < 0, \\
    \text{diverges}        & \text{if } D_I = 0.
  \end{cases}
  \label{eq:sign_flip}
\end{equation}

In practice, the code included a guard $|D_I| < \epsilon \implies
\tilde{\sigma}_I = 0$.  This zeroed out vertex values while midpoint values
(which have $N_{ij} = 4\lambda_i\lambda_j \ge 0$ everywhere) remained correct.
The result was the characteristic checkerboard pattern:

\begin{center}
\begin{tabular}{lcc}
  \toprule
  Node type & Denominator $D_I$ & Projected value \\
  \midrule
  Vertex   & $\le 0$ or $\approx 0$ & $0$ or wrong sign \\
  Midpoint & $> 0$                   & Correct \\
  \bottomrule
\end{tabular}
\end{center}


% ─────────────────────────────────────────────────────────────────────────
\section{Why HEX27 Is Unaffected}
\label{sec:hex27_safe}

HEX27 uses the tensor-product $3 \times 3 \times 3$ Gauss--Legendre rule
(27~points), which has \textbf{all positive weights}.  Furthermore, the
Lagrangian shape functions on the reference cube $[-1,1]^3$ are products of 1D
Lagrange polynomials, and while they can go negative, they do so in a balanced
way that keeps the lumped denominator positive.  The combination of all-positive
weights and the tensor-product structure prevents the denominator catastrophe.


% ─────────────────────────────────────────────────────────────────────────
\section{Resolution: Two Complementary Fixes}
\label{sec:resolution}

Two independent corrections were applied:

\subsection{Fix 1: Switch to an All-Positive-Weight Quadrature Rule}

The root cause is the negative weight $w_1 = -2/15\,|\calS_3|$ in the Keast
5-point rule.  Replacing it with the \textbf{HMS 4-point degree-2 rule}, which
has all positive weights (\cref{sec:HMS}), eliminates the negative-weight
pathway entirely.

In \texttt{SimplexQuadrature.hh}, the default for TET10 was changed:
\begin{lstlisting}
template<>
inline constexpr auto default_simplex_rule<3, 2>() {
    // Was: make_simplex_rule_3d_5()  [Keast 5-pt, w1 = -2/15]
    // Now: make_simplex_rule_3d_4()  [HMS 4-pt, all w > 0]
    return make_simplex_rule_3d_4();
}
\end{lstlisting}

\begin{remark}
The HMS 4-point rule has degree~2, one degree lower than the Keast 5-point
rule (degree~3).  For the TET10 stiffness matrix, degree~$2$ is the minimum
requirement ($2(k-1) = 2$ for $k=2$), so accuracy is maintained for linear
elasticity.  For the mass matrix or high-order nonlinear constitutive models,
the conical product (\texttt{ConicalProductIntegrator<3,2>}, 8~points, degree~3)
is available as a higher-degree drop-in replacement with all positive weights.
\end{remark}

\subsection{Fix 2: Guard Against Negative Denominators}

As a defensive measure, the lumped $L^2$ projection code was modified to use
the absolute value of the denominator:
\begin{lstlisting}
denominator = std::abs(raw_denominator);
if (denominator < epsilon) denominator = epsilon;
\end{lstlisting}

This ensures that even if a future quadrature rule with negative weights is
inadvertently used, the projection produces a \emph{smoothed} (not
sign-flipped) field.  The \texttt{std::abs} guard does not affect the result
when all weights are positive, since the denominator is then guaranteed
positive.

\subsection{Fix 3 (Recommended): SPR Volume-Weighted Averaging}

The most robust solution is to bypass lumped $L^2$ projection entirely and use
\textbf{volume-weighted nodal averaging}\,---\,a simplified form of the
Superconvergent Patch Recovery (SPR) technique~\cite{ZienkiewiczZhu1992a}.
This is discussed in detail in \cref{ch:projection}.


% ─────────────────────────────────────────────────────────────────────────
\section{Summary}

The checkerboard artefact results from the interaction of three factors:

\begin{enumerate}[nosep]
  \item \textbf{Negative quadrature weight} in the Keast 5-point rule
        ($w_1 = -2/15\,|\calS_3|$).
  \item \textbf{Negative vertex shape functions} of TET10 at interior
        quadrature points ($N_i = \lambda_i(2\lambda_i - 1) < 0$ when
        $\lambda_i < 1/2$).
  \item \textbf{Lumped $L^2$ projection} with a denominator that sums signed
        products $N_I \cdot w_g$, allowing cancellation to zero or sign
        reversal at vertex nodes.
\end{enumerate}

Any one of these factors being absent would prevent the artefact.  The chosen
resolution addresses all three:
\begin{itemize}[nosep]
  \item Factor~1: eliminated by switching to HMS 4-point (all $w_g > 0$);
  \item Factor~3: eliminated by switching to SPR averaging (no denominator
        subtraction possible).
  \item Factor~2: is intrinsic to quadratic elements and cannot be changed, but
        it is rendered harmless by the other two fixes.
\end{itemize}
